\chapter{ANTECEDENTES TÉCNICOS: DO SLACGS AO COINFOSIM}
\label{ap:antecedentes}

Este apêndice detalha a evolução técnica mencionada brevemente na introdução (Seção~1.3): a passagem de um simulador antecedente mais simples, o \SLACGS{}, para a formulação atual do \CoInfoSim{}, incluindo a distinção conceitual entre sensores, canais de informação e atributos formais que orienta a terminologia usada em todo o relatório.

\section{O simulador antecedente}

O ponto de partida foi o \SLACGS{}, que utilizava distribuições gaussianas especificadas de forma controlada e simulação de Monte Carlo para acompanhar curvas de perda em diferentes tamanhos amostrais. O \SLACGS{} comparava dimensões por uma \textbf{cadeia incremental prefixada} --- por exemplo $\{X_1\}$, $\{X_1,X_2\}$, $\{X_1,X_2,X_3\}$ --- sem investigar todas as combinações possíveis \cite{azevedo2026coinfosim}. Essa escolha de projeto era adequada ao escopo idealizado do simulador original, mas impedia distinguir, por exemplo, uma dupla formada por atributos fortes mas redundantes de outra composta por atributos individualmente modestos e complementares, e não permitia observar trocas de ranking entre subconjuntos de mesma cardinalidade.

\section{Por que enumerar todos os subconjuntos não vazios}

O \CoInfoSim{} substitui a cadeia incremental prefixada pela enumeração de todos os $2^d-1$ subconjuntos não vazios de atributos. Essa reformulação é o que torna possível: (i) comparar subconjuntos de mesma cardinalidade entre si, não apenas subconjuntos aninhados; (ii) observar diretamente fenômenos de complementaridade e redundância, que dependem da comparação entre configurações alternativas e não apenas do crescimento de uma única cadeia; e (iii) definir o \PredictiveCooperationProfile{} como a evolução do desempenho preditivo relativo entre \emph{todos} os subconjuntos, e não apenas ao longo de uma trajetória de inclusão fixada de antemão. O custo computacional dessa enumeração cresce exponencialmente com $d$ (Capítulo~7), o que motivou manter a dimensionalidade dos cenários deliberadamente baixa (Capítulo~3).

\section{A ancoragem em dados reais}

O \SLACGS{} operava inteiramente sobre modelos gaussianos idealizados, especificados manualmente. A incorporação de datasets públicos exigiu decisões que não apareciam no simulador idealizado: interpretação do domínio, definição de alvos, prevenção de vazamento de informação, particionamento, tratamento de valores ausentes, padronização, balanceamento e análise exploratória da geometria dos dados. O desenvolvimento do projeto passou, assim, a incluir a construção de protocolos científicos específicos para cada conjunto de dados (Capítulo~3), e não apenas a especificação de parâmetros de um modelo gerador.

O protocolo resultante compara três condições experimentais de treinamento, denominadas \emph{braços}: amostras balanceadas retiradas do conjunto de treinamento real; amostras de uma Gaussiana única por classe ajustada apenas ao conjunto de treinamento real; e amostras de uma mistura gaussiana condicional à classe, também ajustada apenas ao conjunto de treinamento real. A avaliação, em todos os braços, ocorre no mesmo conjunto de teste real fixo. Esse desenho se relaciona à lógica \emph{Train on Synthetic, Test on Real}, proposta na literatura de avaliação de dados sintéticos \cite{esteban2017}, mas desloca a ênfase da utilidade preditiva global para a preservação da organização cooperativa que emerge do treinamento dos classificadores para os diferentes subconjuntos.

\section{Sensores, canais de informação e atributos formais}

A ampliação de escopo do \SLACGS{} ao \CoInfoSim{} também motivou uma clarificação terminológica, mantida consistentemente no relatório atual (Capítulo~2): sensores permanecem uma aplicação importante e motivadora, especialmente nos cenários de ocupação e qualidade do ar, mas o SUPPORT2 mostra que a mesma formulação pode representar medidas fisiológicas e laboratoriais, e aplicações futuras podem incluir informações produzidas por pessoas ou por sistemas computacionais. O nome \CoInfoSim{} reflete essa generalização: um \textbf{canal de informação} é a interpretação operacional de uma fonte de aquisição --- sensor, exame, procedimento, registro administrativo ou avaliação humana --- enquanto um \textbf{atributo} é a variável formal que efetivamente alimenta o classificador. Um canal pode corresponder a um único atributo ou a um grupo de atributos explicitamente agrupado; o presente relatório usa \emph{atributo} como termo formal por omissão e reserva \emph{canal de informação} para o sentido operacional aqui descrito, evitando tratar o segundo como sinônimo padrão de toda entrada de um classificador.
